% !TeX root = ../thuthesis-example.tex

\begin{resume}

  % \section*{个人简历}
  \section*{Curriculum Vitae}

  % 197× 年 ×× 月 ×× 日出生于四川××县。
% TO_APPROVE

  Born on XX XX, 197X, in XX County, Sichuan Province.

  % 1992 年 9 月考入××大学化学系××化学专业，1996 年 7 月本科毕业并获得理学学士学位。
% TO_APPROVE

  Admitted to the XX Chemistry program in the Department of Chemistry at XX University in September 1992; graduated in July 1996 with a Bachelor of Science degree.

  % 1996 年 9 月免试进入清华大学化学系攻读××化学博士至今。
% TO_APPROVE

  Admitted without entrance examination to the Department of Chemistry at Tsinghua University in September 1996, and has since been pursuing a doctoral degree in XX Chemistry.


  % \section*{在学期间完成的相关学术成果}
  \section*{Relevant Academic Achievements Completed During Enrollment}

  % \subsection*{学术论文}
  \subsection*{Academic Papers}

  \begin{achievements}
    \item Yang Y, Ren T L, Zhang L T, et al. Miniature microphone with silicon-based ferroelectric thin films[J]. Integrated Ferroelectrics, 2003, 52:229-235.
    % \item 杨轶, 张宁欣, 任天令, 等. 硅基铁电微声学器件中薄膜残余应力的研究[J]. 中国机械工程, 2005, 16(14):1289-1291.
    \item Yang Y, Zhang N X, Ren T L, et al. Study on residual stress in thin films in silicon-based ferroelectric microacoustic devices[J]. China Mechanical Engineering, 2005, 16(14):1289-1291.
    % \item 杨轶, 张宁欣, 任天令, 等. 集成铁电器件中的关键工艺研究[J]. 仪器仪表学报, 2003, 24(S4):192-193.
    \item Yang Y, Zhang N X, Ren T L, et al. Study on key processes in integrated ferroelectric devices[J]. Chinese Journal of Scientific Instrument, 2003, 24(S4):192-193.
    % \item Yang Y, Ren T L, Zhu Y P, et al. PMUTs for handwriting recognition. In press[J]. (已被Integrated Ferroelectrics录用)
    \item Yang Y, Ren T L, Zhu Y P, et al. PMUTs for handwriting recognition. In press[J]. (accepted by Integrated Ferroelectrics)
  \end{achievements}


  % \subsection*{专利}
  \subsection*{Patents}

  \begin{achievements}
    % \item 任天令, 杨轶, 朱一平, 等. 硅基铁电微声学传感器畴极化区域控制和电极连接的方法: 中国, CN1602118A[P]. 2005-03-30.
    \item Ren T L, Yang Y, Zhu Y P, et al. Method for controlling domain polarization regions and electrode connection in a silicon-based ferroelectric microacoustic sensor: China, CN1602118A[P]. 2005-03-30.
    % \item Ren T L, Yang Y, Zhu Y P, et al. Piezoelectric micro acoustic sensor based on ferroelectric materials: USA, No.11/215, 102[P]. (美国发明专利申请号.)
    \item Ren T L, Yang Y, Zhu Y P, et al. Piezoelectric micro acoustic sensor based on ferroelectric materials: USA, No.11/215, 102[P]. (U.S. invention patent application number.)
  \end{achievements}

\end{resume}



% 本科生格式：

% \begin{resume}
%   \section*{学术论文}
%
%   \begin{achievements}
%     \item ZHOU R, HU C, OU T, et al. Intelligent GRU-RIC Position-Loop
%       Feedforward Compensation Control Method with Application to an
%       Ultraprecision Motion Stage[J], IEEE Transactions on Industrial
%       Informatics, 2024, 20(4): 5609-5621.
%
%     \item 杨轶, 张宁欣, 任天令, 等. 硅基铁电微声学器件中薄膜残余应力的研究[J].
%       中国机械工程, 2005, 16(14):1289-1291.
%
%     \item YANG Y, REN T L, ZHU Y P, et al. PMUTs for handwriting recognition.
%       In press[J]. (已被Integrated Ferroelectrics录用)
%
%   \end{achievements}
%
%
%   \section*{专利}
%
%   \begin{achievements}
%     \item 胡楚雄, 付宏, 朱煜, 等. 一种磁悬浮平面电机: ZL202011322520.6[P]. 2022-04-01.
%
%     \item REN T L, YANG Y, ZHU Y P, et al. Piezoelectric micro acoustic sensor
%       based on ferroelectric materials: No.11/215, 102[P]. (美国发明专利申请号.)
%
%   \end{achievements}
% \end{resume}
